\documentclass[12pt, a4paper]{article}
\usepackage[utf8]{inputenc}
\usepackage[spanish,es-noshorthands]{babel}
\usepackage[margin=1.5cm]{geometry}
\usepackage{tabularx}
\usepackage{multirow}
\usepackage{array}
\usepackage{amssymb}
\usepackage[table]{xcolor}

\begin{document}

\begin{center}
    \textbf{UNIVERSIDAD DE ORIENTE}\\
    \textbf{VICERRECTORADO ACADÉMICO}\\
    \textbf{CONSEJO DE INVESTIGACIÓN}\\
    \vspace{0.5cm}
    {\Large \textbf{INSCRIPCIÓN / ACTUALIZACIÓN DE}}\\
    {\Large \textbf{GRUPOS DE INVESTIGACIÓN}}
\end{center}

\vspace{0.5cm}

\noindent
\renewcommand{\arraystretch}{1.5}
\begin{tabularx}{\textwidth}{|X|X|X|}
\hline
\multicolumn{3}{|c|}{\cellcolor[gray]{0.9}\textbf{IDENTIFICACIÓN DEL INVESTIGADOR (Coordinador del Grupo)}} \\
\hline
\multicolumn{2}{|l|}{Rodríguez Gómez} & César Augusto \\
\hline
C.I. N$^\circ$ 5701716& \multicolumn{2}{l|}{Lic. en Física} \\
\hline
\multicolumn{2}{|l|}{Departamento de Física} & Escuela de Ciencias \\
\hline
8 horas/dia dedicadas a la investigación & Asistente en investigación & Dedicación Exclusiva \\
\hline
Telf. Hab & Celular: (+58)412-188.31.84 & Telf. Ofic. \\
\hline
E-Mail: cero97@yahoo.com & \multicolumn{2}{l|}{Av. Cancamure, Urb. Los Angeles, Casa 27} \\
\hline
\multicolumn{2}{|l|}{Ultima premiación obtenida} & Período \\
\hline
\multicolumn{1}{|l|}{Tipo de Investigador} & \multicolumn{2}{l|}{Indique:} \\
\multicolumn{1}{|l|}{\quad $\blacksquare$ Novel} & \multicolumn{2}{l|}{\quad 0009-0007-2674-2466 - No. ORCID} \\
\multicolumn{1}{|l|}{\quad $\square$ Con Publicación (CP)} & \multicolumn{2}{l|}{\quad \rule{4cm}{0.4pt} Índice h} \\
\multicolumn{1}{|l|}{\quad $\square$ Investigador Activo} & \multicolumn{2}{l|}{\quad \rule{4cm}{0.4pt} Índice i-10} \\
\multicolumn{1}{|l|}{} & \multicolumn{2}{l|}{\quad \rule{4cm}{0.4pt} Otros índices*} \\
\multicolumn{1}{|l|}{} & \multicolumn{2}{l|}{*Explique \rule{6cm}{0.4pt}} \\
\hline
\end{tabularx}

\vspace{0.5cm}

\noindent
\begin{tabularx}{\textwidth}{|X|X|}
\hline
\multicolumn{2}{|c|}{\cellcolor[gray]{0.9}\textbf{IDENTIFICACIÓN DEL GRUPO DE INVESTIGACIÓN}} \\
\hline
\multicolumn{2}{|l|}{Nombre del Grupo: GIDEAL (Grupo de Invest. en Diseño Electrónico y Aprendizaje con LLMs)} \\
\hline
\multirow{2}{=}{\textbf{Objetivos:} {\footnotesize Desarrollar sistemas basados en agentes IA y LLMs para asistir en la enseñanza, el diseño electrónico (EDA) y la evaluación académica.}} & \textbf{Palabras Claves:} {\footnotesize IA, LLMs, EDA, RAG, Educación, Agentes} \\
\cline{2-2}
 & \textbf{Área de Investigación:} {\footnotesize Inteligencia Computacional Aplicada} \\
\hline
\multicolumn{2}{|p{\dimexpr\textwidth-2\tabcolsep-2\arrayrulewidth\relax}|}{\textbf{Descripción del Grupo:} \newline {\footnotesize GIDEAL es un equipo multidisciplinario enfocado en la convergencia de la Ingeniería Electrónica e Inteligencia Artificial. Diseñamos flujos orquestados por IA bajo una arquitectura de tres capas. Nuestras actividades incluyen la creación de asistentes conversacionales (RAG), extracción automatizada de esquemas EDA desde LaTeX, evaluación multimodal de prácticas con visión artificial y desarrollo de servidores MCP con pasarela de servicios de mensajería por internet (Telegram, WhatsApp, Signal, Discord), con el fin de optimizar el trabajo académico y docente.}} \\[1.5cm]
\hline
\end{tabularx}

\newpage

\noindent
\begin{tabularx}{\textwidth}{|X|X|X|}
\hline
\multicolumn{3}{|c|}{\cellcolor[gray]{0.9}\textbf{IDENTIFICACIÓN DE LOS INTEGRANTES}} \\
\hline
Investigador & Escuela/Departamento & Firmas \\
\hline
César Rodríguez & Ciencias/Física & \\ \hline
 & & \\ \hline
 & & \\ \hline
 & & \\ \hline
 & & \\ \hline
 & & \\ \hline
 & & \\ \hline
 & & \\ \hline
 & & \\ \hline
 & & \\ \hline
\end{tabularx}

\newpage

\noindent
\begin{tabularx}{\textwidth}{|X|}
\hline
\multicolumn{1}{|c|}{\cellcolor[gray]{0.9}\textbf{ACTIVIDADES DEL GRUPO DE INVESTIGACIÓN}} \\
\hline
\vspace{0.2cm}
\begin{itemize}\small
    \item \textbf{Asistencia Educativa Automatizada:} Desarrollo e integración de chatbots académicos (sistemas RAG) para procesamiento y consulta de material de estudio, exámenes y prácticas en formatos PDF, Markdown y LaTeX.
    \item \textbf{Análisis Multimodal de Prácticas y Exámenes:} Orquestación de flujos de trabajo con LLMs multimodales (como Gemini) para realizar evaluaciones preliminares (sujetas a revisión del profesor), análisis y diagnóstico de exámenes escritos e informes de laboratorio.
    \item \textbf{Desarrollo de Agentes EDA:} Creación de agentes inteligentes capaces de extraer netlists y listas de materiales (BOM) a partir de esquemas en LaTeX (circuitikz) para su integración en herramientas de diseño electrónico (como EasyEDA o KiCAD).
    \item \textbf{Automatización y Orquestación Remota:} Implementación y gestión de servidores locales MCP y pasarelas (gateways) a través de servicios de mensajería para acceder y automatizar flujos de investigación sin requerir configuraciones de red complejas.
    \item \textbf{Creación de Recursos Docentes:} Diseño automatizado de rúbricas, guías y material instruccional destinado a las asignaturas de Física y Electrónica mediante inteligencia artificial generativa.
\end{itemize}
\vspace{10cm} \\
\hline
\end{tabularx}

\newpage

\noindent
\begin{tabularx}{\textwidth}{|X|}
\hline
\textbf{Equipos y/o Materiales del Grupo:}\vspace{0.2cm}
\begin{itemize}\small
    \item \textbf{Hardware de Cómputo:} Computadores y estaciones de trabajo para la ejecución concurrente de agentes de IA, bases de datos vectoriales (ChromaDB) y modelos de lenguaje locales o de código abierto (vLLM).
    \item \textbf{Instrumentación Electrónica:} Osciloscopios, generadores de funciones, fuentes de alimentación y multímetros digitales para las validaciones prácticas y desarrollo de dispositivos (IoT, fuentes SMPS).
    \item \textbf{Servidores y Pasarelas:} Servidores locales configurados con servicios de pasarela remota (Telegram Gateway, Waydroid) para la orquestación y diagnóstico del flujo de automatización docente.
    \item \textbf{Herramientas de Desarrollo:} Acceso a plataformas EDA (KiCAD, EasyEDA) y créditos/APIs para grandes modelos de lenguaje multimodales (Groq, Gemini, OpenRouter).
    \item \textbf{Componentes y Materiales de Laboratorio:} Componentes electrónicos discretos (BJT, MOSFET, condensadores), placas de desarrollo para prototipado y material fungible para el desarrollo de prácticas físicas.
\end{itemize}
\vspace{6cm} \\
\hline
Favor anexar\\
\begin{itemize}
    \item Carta Presentación del Grupo
    \item Curriculum actualizado de cada uno de los integrantes
    \item Lista de Proyectos Realizados y en Ejecución
    \item Lista y copia(s) de (los) Artículos Publicados en los últimos 2 años
    \item Lista de trabajos de pre-grado y post-grado en los últimos 2 años
\end{itemize}\\
\hline
\end{tabularx}

\vspace{0.5cm}

\noindent
\begin{tabularx}{\textwidth}{|X|X|}
\hline
\multicolumn{2}{|c|}{\cellcolor[gray]{0.9}\textbf{FORMALIZACIÓN DE INSCRIPCIÓN/ACTUALIZACIÓN}} \\
\hline
Coordinador del Grupo & Registrado y Conformado por el Consejo de Investigación \\
\hline
 & \\
 & \\
 & \\
Firma y fecha & Firma y fecha \\
\hline
\end{tabularx}

\end{document}
